\subsection{GIT 分析}\label{src:9.5.3-subsection}
目标是证明定理~\bookref{9.5.1}: 存在整数 $D$ 和 $M$, 使对所有 $d>D$, $m>M$, 商 $[\mathcal O_C\otimes V\twoheadrightarrow F]\in\Quot^P(\mathcal O_C\otimes V/C)$ 关于 $L_m$ 和 $\mathrm{SL}(V)$ 的作用是 GIT 半稳定的, 当且仅当 $F$ 是半稳定向量丛, 且 $\mathrm H^0(C,\mathcal O_C\otimes V)\to\mathrm H^0(C,F)$ 为同构.
\begin{proposition}[GIT 半稳定性判据]\label{src:9.5.12}
设 $P$ 是秩 $r$, 次数 $d$ 的向量丛的 Hilbert 多项式, 置 $V=\Bbbk^{\oplus P(0)}$. 对每个秩 $r>0$ 和次数 $d$, 存在整数 $M$, 使对所有 $m>M$, 对商 $[\mathcal O_C\otimes V\twoheadrightarrow F]\in\Quot^P(\mathcal O_C\otimes V/C)$, 下列条件等价:
\begin{enumerate}[(1)]
\item $[\mathcal O_C\otimes V\twoheadrightarrow F]$ 关于 $L_m$ 和 $\mathrm{SL}(V)$ 是 GIT 半稳定的;
\item 对每个子空间 $V'\subseteq V$, 令 $F':=\im(\mathcal O_C\otimes V'\to F)$, 有
\begin{equation}\label{src:9.5.13}
\frac{\chi(C,F'(m))}{\dim V'}\geq\frac{\chi(C,F(m))}{\dim V};
\end{equation}
\item[(2$'$)] 对每个子空间 $V'\subseteq V$, 令 $F':=\im(\mathcal O_C\otimes V'\to F)$, 有
\[
\frac{P_{F'}(z)}{\dim V'}\geq\frac{P(z)}{\dim V},
\]
其中 $P_{F'}(z)$ 为 $F'$ 的 Hilbert 多项式.
\end{enumerate}
\end{proposition}
\begin{proof}
设 $\lambda:\mathbb G_m\to\mathrm{SL}(V)$ 为单参数子群, 将 $V=\bigoplus_{n\in\mathbb Z}V_n$ 写为权空间的直和, 其中 $V_n$ 的权重为 $n$. 因为 $\lambda$ 的像在 $\mathrm{SL}(V)$ 中, 有 $\sum_n n\dim V_n=0$. 置 $V_{\leq n}:=\bigoplus_{j\leq n}V_j$, 得到滤过
\[
V_\bullet:0\subseteq\cdots\subseteq V_{\leq n-1}\subseteq V_{\leq n}\subseteq\cdots\subseteq V.
\]
令 $F_{\leq n}$ 为 $\mathcal O_C\otimes V_{\leq n}\hookrightarrow\mathcal O_C\otimes V\twoheadrightarrow F$ 的像, 得到诱导滤过
\[
F_\bullet:0\subseteq\cdots\subseteq F_{\leq n-1}\subseteq F_{\leq n}\subseteq\cdots\subseteq F.
\]
定义 $F_n:=F_{\leq n}/F_{\leq n-1}$, 有满射 $\mathcal O_C\otimes V_n\to F_n$, 从而有满射 $[\mathcal O_C\otimes V\twoheadrightarrow\mathrm{gr}(F_\bullet)]\in\Quot^P(\mathcal O_C\otimes V/C)$, 其中 $\mathrm{gr}(F_\bullet):=\bigoplus_n F_n$. 类似命题~\bookref{7.9.6} 的论证, 可验证 $\lim_{t\to0}\lambda(t)\cdot[F]=\mathrm{gr}(F_\bullet)$. 因此
\[
\mu_{L_m}([F],\lambda)=-\sum_{n\in\mathbb Z}n\chi(C,F_n(m)).
\]
利用 $\sum_n n\dim V_n=0$, 将 Hilbert--Mumford 指标改写为
\[
\begin{aligned}
\mu_{L_m}([F],\lambda)&=-\sum_{n\in\mathbb Z}n\chi(C,F_n(m))+\frac{\chi(C,F(m))}{\dim V}\sum_{n\in\mathbb Z}n\dim V_n\\
&=-\frac1{\dim V}\sum_{n\in\mathbb Z}n\big(\dim V\chi(C,F_n(m))-\dim V_n\chi(C,F(m))\big)\\
&=-\frac1{\dim V}\sum_{n\in\mathbb Z}n\Big(\dim V\big(\chi(C,F_{\leq n}(m))-\chi(C,F_{\leq n-1}(m))\big)\\
&\hspace{65pt}-(\dim V_{\leq n}-\dim V_{\leq n-1})\chi(C,F(m))\Big)\\
&=\frac1{\dim V}\sum_{n\in\mathbb Z}\big(\dim V\chi(C,F_{\leq n}(m))-\dim V_{\leq n}\chi(C,F(m))\big),
\end{aligned}
\]
其中每个求和只有有限个非零项.

为说明 (2) $\Rightarrow$ (1), 对每个 $n$, 将不等式~\bookeqref{9.5.13} 应用于 $V'=V_{\leq n}$. 得到 $\dim V\chi(C,F_{\leq n}(m))\geq\dim V_{\leq n}\chi(C,F(m))$, 从而 $\mu_{L_m}([F],\lambda)\geq0$. 反之, 假设存在违反~\bookeqref{9.5.13} 的子空间 $V'\subseteq V$. 选取补空间 $V''\subseteq V$, 使 $V\cong V'\oplus V''$, 并选取 $V'$ 与 $V''$ 的基. 设 $\lambda:\mathbb G_m\to\mathrm{SL}(V)$ 为关于所选基的对角单参数子群, 在 $V'$ 上权重为 $-\dim V''$, 在 $V''$ 上权重为 $\dim V'$. 则
\[
\mu([F],\lambda)=\dim V\chi(C,F'(m))-\dim V'\chi(C,F(m))<0.
\]
为说明 (2) $\Leftrightarrow$ (2$'$), 注意当 $V'\subseteq V$ 遍历所有子空间时, 层族 $F'=\im(\mathcal O_C\otimes V'\to F)$ 有界. 故 Hilbert 多项式 $P_{F'}$ 只有有限种选择. 因为每个 $P_{F'}$ 都是一次的, 可增大 $M$, 使对每种可能的 $P_{F'}$, 当 $m>M$ 时有
\[
\frac{P_{F'}(z)}{\dim V'}\geq\frac{P_F(z)}{\dim V}\iff\frac{P_{F'}(m)}{\dim V'}\geq\frac{P_F(m)}{\dim V}.
\]
另见 \cite[引理 1.15--16]{Sim94a}, \cite[引理 7.2.2]{Lan75} 及 \cite[引理 4.4.5--6]{HL10}.
\end{proof}
\begin{exercise}[不重要]\label{src:9.5.14}
将命题~\bookref{9.5.12} 的等价条件扩充为包含:
\begin{enumerate}[(1)]
\item[(3)] 对每个子层 $F'\subseteq F$, 令 $V':=\mathrm H^0(C,F'(m))$, 有 $\chi(C,F'(m))/\dim V'\geq\chi(C,F(m))/\dim V$.
\end{enumerate}
\end{exercise}
\begin{proof}[定理~\bookref{9.5.1} 的证明]
选取 $D>r(2g-1)$, 使半稳定性判据 (\bookref{9.5.5}) 成立, 再选取 $M$, 使 GIT 半稳定性判据 (\bookref{9.5.12}) 成立. 设 $[\mathcal O_C\otimes V\twoheadrightarrow F]\in\Quot^P(\mathcal O_C\otimes V/C)$ 为商. 需要证明当 $d>D$, $m>M$ 时,
\[
[\mathcal O_C\otimes V\twoheadrightarrow F]\text{ 关于 }L_m\text{ GIT 半稳定}\iff F\text{ 半稳定且 }V\xrightarrow{\sim}\mathrm H^0(C,F).
\]
对 ($\Leftarrow$), 设 $V'\subseteq V$ 为子空间, $F'=\im(\mathcal O_C\otimes V'\to F)$. 交换图
\[
\begin{tikzcd}[ampersand replacement=\&]
V'\arrow[r]\arrow[d,hook]\&\mathrm H^0(C,F')\arrow[d,hook]\\
V\arrow[r,"\sim"]\&\mathrm H^0(C,F)
\end{tikzcd}
\]
说明 $\dim V'\leq\mathrm h^0(F')$. 因为 $F$ 半稳定, 半稳定性判据 (\bookref{9.5.5}) 给出 $\mathrm h^0(F')/\mathrm{rk}\,F'\leq\mathrm h^0(F)/\mathrm{rk}\,F$. 故
\begin{equation}\label{src:9.5.15}
\frac{\dim V'}{\mathrm{rk}\,F'}\leq\frac{\mathrm h^0(F')}{\mathrm{rk}\,F'}\leq\frac{\mathrm h^0(F)}{\mathrm{rk}\,F}=\frac{\dim V}{\mathrm{rk}\,F},
\end{equation}
可改写为
\begin{equation}\label{src:9.5.16}
\frac{\mathrm{rk}\,F'}{\dim V'}\geq\frac{\mathrm{rk}\,F}{\dim V}.
\end{equation}
根据 GIT 半稳定性判据 (\bookref{9.5.12}), 只需证明
\begin{equation}\label{src:9.5.17}
\frac{P_{F'}(z)}{\dim V'}\geq\frac{P_F(z)}{\dim V}.
\end{equation}
两边首项由~\bookeqref{9.5.16} 中的表达式给出, 故若~\bookeqref{9.5.16} 中严格不等式成立, 此不等式也成立. 若~\bookeqref{9.5.16} 中等号成立, 则~\bookeqref{9.5.15} 中各项均相等, 半稳定性判据 (命题~\bookref{9.5.5}) 蕴含 $\mu(F')=\mu(F)$. 因此~\bookeqref{9.5.17} 的常数项满足
\[
\frac{d'+r'(1-g)}{\dim V'}=\frac{r'(\mu(F')+1-g)}{\dim V'}\geq\frac{r(\mu(F)+1-g)}{\dim V}=\frac{d+r(1-g)}{\dim V},
\]
其中 $r'=\mathrm{rk}\,F'$, $d'=\deg F'$. 故~\bookeqref{9.5.17} 成立.

反之, 对 ($\Rightarrow$), 假设 $[\mathcal O_C\otimes V\twoheadrightarrow F]$ GIT 半稳定. 先证明 $\mathrm H^1(C,F)=0$. 由 Serre 对偶 (\bookref{6.1.3}), $\mathrm H^1(C,F)\cong\mathrm{Hom}_C(F,\Omega_C)^\vee$. 假设 $F\to\Omega_C$ 是非零映射. 令 $F''=\im(F\to\Omega_C)$, $V''=\mathrm H^0(C,F'')$, $F'=\ker(F\to F'')$, $V'=\mathrm H^0(C,F')=\ker(V\to V'')$. 则 $P_F(z)=P_{F'}(z)+P_{F''}(z)$, $\dim V\leq\dim V'+\dim V''$. $I:=\im(\mathcal O_C\otimes V'\to F)$ 包含于 $F'$, 故 $P_{F'}(z)\geq P_I(z)$. GIT 半稳定性判据 (\bookref{9.5.12}) 表明
\[
\frac{P_{F'}(z)}{\dim V'}\geq\frac{P_I(z)}{\dim V'}\geq\frac{P_F(z)}{\dim V},
\]
不难看出这蕴含商 $F\twoheadrightarrow F''$ 满足
\begin{equation}\label{src:9.5.18}
\frac{P_{F''}(z)}{\dim V''}\leq\frac{P_F(z)}{\dim V}.
\end{equation}
比较首项得 $\mathrm{rk}\,F''/\dim V''\leq\mathrm{rk}\,F/\dim V$. 对此不等式取倒数, 并代入 $\dim V=r(1-g)+d$, $\dim V''\leq\dim\mathrm H^0(C,\Omega_C)=g$, $\mathrm{rk}\,F''=1$ (因为假设 $F''\subseteq\Omega_C$ 非零!), 得
\[
g=\frac{g}{\mathrm{rk}\,F''}\geq\frac{\dim V''}{\mathrm{rk}\,F''}\geq\frac{\dim V}{\mathrm{rk}\,F}=\mu(F)+1-g.
\]
这蕴含 $\mu(F)\leq2g-1$, 与 $d>D>r(2g-1)$ 矛盾. 故 $\mathrm H^1(C,F)=0$.

其次, 证明自然映射 $V\to\mathrm H^0(C,F)$ 为同构. 因为 $\mathrm H^1(C,F)=0$, 两个空间维数相同, 只需证明映射单射. 若有非零核 $V'=\ker(V\to\mathrm H^0(C,F))$, 则 $F'=\im(\mathcal O_C\otimes V'\to F)=0$, $P_{F'}(z)=0$, 违反 GIT 半稳定性判据 (\bookref{9.5.12}).

最后证明 $F$ 半稳定. 根据半稳定性判据 (\bookref{9.5.5}), 只需证明对所有子层 $0\neq E\subseteq F$, 有 $\mathrm h^0(E)/\mathrm{rk}\,E\leq\mathrm h^0(F)/\mathrm{rk}\,F$. 可假设 $\mathrm h^0(E)>0$. 令 $V'=\mathrm H^0(C,E)\subset V$, 令 $F'\subseteq E\subset F$ 为 $\mathcal O_C\otimes V'\to E$ 的像. GIT 半稳定性判据 (\bookref{9.5.12}) 蕴含 $P_{F'}(z)/\dim V'\geq P_F(z)/\dim V$, 比较首项得
\[
\frac{\mathrm{rk}\,E}{\mathrm h^0(E)}=\frac{\mathrm{rk}\,E}{\dim V'}\geq\frac{\mathrm{rk}\,F'}{\dim V'}\geq\frac{\mathrm{rk}\,F}{\dim V}=\frac{\mathrm{rk}\,F}{\mathrm h^0(F)},
\]
即所需不等式的乘法倒数.

另见 \cite[定理 1.19]{Sim94a}, \cite[引理 7.2.4]{LP97} 及 \cite[定理 4.3.3]{HL10}.
\end{proof}
\begin{exercise}\label{src:9.5.19}
陈述并证明定理~\bookref{9.5.1} 关于稳定向量丛的版本.
\end{exercise}
\begin{remark}[箭图构造]\label{src:9.5.20}
\'Alvarez-C\'onsul 与 King 在 \cite{CK07} 中给出了 $M_{r,d}^{\mathrm{ps}}$ 的另一种构造, 也适用于高维簇上层的模空间. 如同上述 GIT 方法, 当 $n\gg0$ 时, 满射 $\mathrm H^0(C,F(n))\otimes\mathcal O_C(-n)\twoheadrightarrow F$ 对 $m-n\gg0$ 诱导满射
\begin{equation}\label{src:9.5.21}
\mathrm H^0(C,F(n))\otimes\mathrm H^0(C,\mathcal O_C(m-n))\twoheadrightarrow\mathrm H^0(C,F(m)).
\end{equation}
我们将~\bookeqref{9.5.21} 看作 Kronecker 箭图的表示, 而不选取同构 $V\xrightarrow{\sim}\mathrm H^0(C,F)$. 此箭图有两个顶点, 从第一个顶点到第二个顶点有 $\mathrm h^0(\mathcal O_C(m-n))$ 条箭头, 表示的维数分别为 $\mathrm h^0(F(n))$ 与 $\mathrm h^0(F(m))$. 对箭图表示, 其 GIT 稳定性有简单的子表示解释 (见习题~\bookref{8.5.6}). 作者随后大致证明 $F$ 半稳定当且仅当箭图表示~\bookeqref{9.5.21} GIT 半稳定. 然而, 如本节开头引文所述, 要完成构造仍须解决额外的技术细节.
\end{remark}
